% B.1--B.3. Standard amsmath/amssymb/amsthm commands only.
\section{B.1. Counting pinned maps without a rigidity assumption}

We use left actions and compose functions from right to left. An oriented
combinatorial map is a triple $C=(E,\rho,\alpha)$, where $E$ is a finite
set of darts, $\rho$ is a permutation of $E$, and $\alpha^2=1$.
We put $\phi=\rho^{-1}\alpha$. Vertices, edges and faces are, respectively,
the cycles of $\rho,\alpha,\phi$, with their cycle lengths as degrees.
An edge of degree one is a half-edge. The map is clean when $\alpha$
has no fixed points. Connectedness means that $\langle\rho,\alpha\rangle$
is transitive. An isomorphism $C\longrightarrow C'$ is a bijection
$u:E\longrightarrow E'$ commuting with the indicated permutations.
These conventions agree with $\rho_f=\rho_s^{-1}\sigma$ on cote 138,
scan 106. Replacing vertex rotations by their inverses changes the
orientation convention, but none of the counts or fields below.

\paragraph{B0.1 (Known: permutations, coverings and dessins).}
Finite transitive permutation data of this kind describe connected
pre-clean Belyi pairs, up to isomorphism over the target
$\mathbb P^1_{\mathbb C}$; the cycle lengths are ramification indices.
For a clean pair $(X,\beta)$, the underlying ordinary map has vertices
$\beta^{-1}(0)$, edges the components of
$\beta^{-1}([0,1])\setminus\beta^{-1}(0)$, and faces the complementary
disks. The Belyi automorphism group identifies with
$\operatorname{Cent}_{\operatorname{Sym}(E)}(\rho,\alpha)$.
We use Schneps, Lemma I.1 and Theorem I.5, with Definition 1 for
pre-clean versus clean, and Jones, Theorem 1 and Sections 2--3
\cite{Schneps1994B,Jones2018B}. These sources have been opened;
the Jones arXiv \LaTeX{} source has also been read.
For the genus convention, the Euler formula is
$\#\operatorname{cyc}(\rho)-\#\operatorname{cyc}(\alpha)
+\#\operatorname{cyc}(\phi)=2-2g$ in the clean case.
For a general dessin of degree $d$, it is
$\#\operatorname{cyc}(\rho)+\#\operatorname{cyc}(\alpha)
+\#\operatorname{cyc}(\phi)=d+2-2g$.

Fix a finite precarte
$T=(S,A,F;p,q,r)$ with $q\in\{1,2\}$ as on scans 98--99.
Write $a_n,b_n,c_n$ for the numbers of its vertices, edges and faces of
degree $n$, and set
\[
 \Gamma_T=\operatorname{Aut}(T),\qquad
 \gamma_T=|\Gamma_T|=\prod_{n\geq1}a_n!b_n!c_n!.
\]
The functor $\operatorname{Fac}$ forgets incidences and cyclic order,
retaining these three graded sets. A $T$-pinning of $C$ is an isomorphism
$u:T\longrightarrow\operatorname{Fac}(C)$ in the category of precartes.
We count isomorphism classes, not objects with a chosen representative
dart set. Let $\mathcal M(T)$ be the finite set of unpinned isomorphism
classes of the prescribed type, with any explicitly imposed connectedness
or genus restriction. For a representative $C$, put
\[
 A_C=\operatorname{Aut}(C),\quad
 K_C=\ker\bigl(A_C\longrightarrow
              \operatorname{Aut}(\operatorname{Fac}(C))\bigr),\quad
 I_C=\operatorname{im}\bigl(A_C\longrightarrow
              \operatorname{Aut}(\operatorname{Fac}(C))\bigr).
\]

\paragraph{B1.1 (Proved: the exact orbit-size formula).}
The number of isomorphism classes of $T$-pinned oriented maps is
\begin{equation}\label{eq:B11}
 N_T=\sum_{[C]\in\mathcal M(T)}\frac{\gamma_T}{|I_C|}
     =\gamma_T\sum_{[C]\in\mathcal M(T)}\frac{|K_C|}{|A_C|}.
\end{equation}
The automorphism group of every pinning of $C$ is $K_C$.
The manuscript's expression $|\mathcal M(T)|\gamma_T$ is correct
if and only if every $I_C$ is trivial. Rigidity is instead the condition
that every $K_C$ is trivial; these are different conditions.

\begin{proof}
For a fixed $C$, the set
$\operatorname{Iso}(T,\operatorname{Fac}(C))$ is a right $\Gamma_T$-torsor
under precomposition, hence has $\gamma_T$ elements. The group $A_C$
acts on its left by postcomposition with $\operatorname{Fac}(a)$.
Two pinnings give isomorphic pinned maps precisely when they lie in
one such orbit: this is exactly the commutative square defining a pinned
isomorphism. Since $u$ is bijective,
$\operatorname{Fac}(a)u=u$ holds exactly when $a\in K_C$.
Every orbit therefore has size $|A_C|/|K_C|=|I_C|$. Dividing the
torsor cardinality by this size and summing over distinct unpinned
classes proves the formula. Every summand is at most $\gamma_T$,
and equality occurs exactly when $|I_C|=1$; positivity proves the stated
criterion for the whole sum, including the vacuous empty case.
Nothing in this proof uses connectedness, so it applies to disconnected
finite maps as well.
\end{proof}

There are no maps unless the three total degrees are equal. This necessary
condition alone need not guarantee a map with additional restrictions.
When $d=0$, our permutation category contains exactly the empty map;
it contributes one if empty maps are allowed, and contributes nothing
if connected nonempty maps are required. When $d=1$, the only datum is
$\rho=\alpha=1$ and there is one pre-clean map and one pinning.
There are no nonempty clean maps of odd dart degree. These conventions
also specify the small-degree cases of (\ref{eq:B11}).

\paragraph{B1.2 (Refuted: the rigidity assertion on 138:99).}
Let $E=\mathbb Z/4$, $\rho(x)=x+1$, and $\alpha(x)=x+2$.
Then $\phi(x)=x+1$; the map has one vertex, two edges, one face,
and genus one. Its full combinatorial automorphism group is $C_4$,
its pinned automorphism group is $C_2$, and its type has one pinned
isomorphism class rather than the two predicted by the manuscript.

\begin{proof}
The permutation $\rho$ is transitive, $\alpha$ is fixed-point-free of
order two, and $\phi=\rho^{-1}\alpha=\rho$. Thus all defining hypotheses
hold. Every permutation commuting with the four-cycle $\rho$ is a
power of $\rho$, and every such power commutes with $\alpha$; hence
$A_C=\langle\rho\rangle$. Its even powers fix both edge sets
$\{0,2\},\{1,3\}$, while its odd powers exchange them. The vertex and
face are unique. Consequently $K_C=\{1,\rho^2\}$ and $|I_C|=2$.
The type has $\Gamma_T\simeq S_2$, so its contribution is one.
There is only one unpinned map of this type: after making $\rho$ the
standard four-cycle, the three fixed-point-free involutions are
$(02)(13)$, $(01)(23)$ and $(03)(12)$; only the first gives one face.
The certificate enumerates all 24 relabellings and verifies all groups
and cycles, not merely the existence of the exhibited involution.
\end{proof}

This is the same finite counterexample recorded in the attached study
and credited there to the Pro notes and Michel Hua's project. We have
re-derived it rather than importing their verification. A rigid example
already shows why rigidity would not justify the manuscript's count:
the one-loop sphere has a $C_2$ automorphism exchanging its two faces,
so $K_C=1$ while $|I_C|=2$.

\paragraph{B1.3 (Proved: all connected cellwise combinatorial kernels).}
For a connected nonempty finite map, $K_C$ is nontrivial if and only if,
for some integer $g\geq1$, its dart permutations are isomorphic to
\[
 E=\mathbb Z/(4g),\qquad \rho(x)=x+1,\qquad \alpha(x)=x+2g.
\]
In this case $K_C=\{1,\alpha\}\simeq C_2$, and the map has one vertex,
$2g$ edges and one face on a genus-$g$ surface. In particular, a half-edge
forces $K_C=1$.

\begin{proof}
An automorphism of a transitive permutation representation fixing one
dart fixes every dart: if $a(x)=x$, then
$a(w x)=w a(x)=w x$ for every word $w$ in $\rho,\alpha$.
Thus $A_C$ acts freely on $E$. If $1\ne a\in K_C$, preservation of
each $\alpha$-cycle gives $a(x)\in\{x,\alpha x\}$, and freeness gives
$a(x)=\alpha x$ for every $x$. There can be no fixed point of $\alpha$,
and $a=\alpha$ is the only possible nonidentity kernel element.
It commutes with $\rho$ and preserves each $\rho$-cycle. Consequently
$\langle\rho,\alpha\rangle$ preserves each $\rho$-cycle, so transitivity
forces $\rho$ to be a single cycle. Write $|E|=2m$.
The unique nontrivial involution in the centralizer of this cycle is
$\alpha=\rho^m$. Now $\phi=\rho^{m-1}$ has
$\gcd(2m,m-1)=\gcd(2,m-1)$ cycles. Translation by $m$ preserves each
such cycle if and only if this gcd divides $m$, which holds exactly when
$m$ is even. Writing $m=2g$ gives the asserted model, with a single
face. Its genus follows from $1-2g+1=2-2g$. Conversely all stated
permutations commute and $\alpha$ fixes the sole vertex, every edge,
and the sole face, proving sufficiency.
\end{proof}

\section{B.2. The three sextic trees and their fields}

We now consider bipartite dessins of degree six and ordered passport
$[3,2,1;2,2,1,1;6]$, with black vertices over $0$, white vertices over
$1$, and the face over $\infty$. These are pre-clean maps with two
half-edges, or equivalently bipartite plane trees with six edges.
We keep the target points $0,1,\infty$ fixed throughout; target
permutations are not part of the isomorphism relation.

\paragraph{B2.1 (Proved: the completed remainder calculation, 138:107--112).}
Define polynomials over $\mathbb Q$ by
\begin{align*}
 h_c(z)&=z^3(z-1)^2(z-c),&
 G_c(z)&=z^2-\frac{5c+4}{6}z+\frac c2,\\
 \delta(c)&=25c^2-32c+16,&
 p(c)&=25c^3-12c^2-24c-16,\\
 P(c)&=-\frac{\delta(c)p(c)}{7776},&
 Q(c)&=\frac{c(5c-8)(25c^3-6c^2+8)}{2592}.
\end{align*}
Then $h'_c(z)=6z^2(z-1)G_c(z)$, and the remainder of $h_c(z)$
on division by $G_c(z)$ is $P(c)z+Q(c)$. In particular, the literal
coefficient called $P(c)$ on scan 109 is quintic, not cubic.
After requiring distinct residual critical points, its effective
equation is the cubic $p(c)=0$. Modulo $p(c)$,
\begin{align}
 Q(c)&=\frac{32-52c-c^2}{1500},\label{eq:B2Q}\\
 h_c(z)-Q(c)&=G_c(z)^2H_c(z),\label{eq:B2factor}\\
 H_c(z)&=z^2+\frac{2c-2}{3}z+\frac{5c^2-4c-4}{12}.\nonumber
\end{align}

\begin{proof}
Differentiating the factored polynomial gives the displayed derivative.
In the quotient $\mathbb Q[c,z]/(G_c)$ we have
$z^2=\alpha z+\beta$, with
$\alpha=(5c+4)/6$ and $\beta=-c/2$.
Set
\[
 A=(\alpha-2)(\alpha-c)+\beta+1,
 \qquad B=(\alpha-2)\beta-(\beta+1)c.
\]
In the quotient, $(z-1)^2(z-c)=Az+B$ and
$z^3=(\alpha^2+\beta)z+\alpha\beta$. Their product gives
$P=A(\alpha^3+2\alpha\beta)+(\alpha^2+\beta)B$ and
$Q=\beta(\alpha^2+\beta)A+\alpha\beta B$.
These formulas decide both the $(\alpha-c)$ versus $\alpha$ reading
and the minus sign preceding $(\beta+1)c$ on scan 108.
Its coefficients,
written without factorizations, are
\begin{align*}
 P(c)&=\frac{-625c^5+1100c^4-184c^3-176c^2-128c+256}{7776},\\
 Q(c)&=\frac{125c^5-230c^4+48c^3+40c^2-64c}{2592}.
\end{align*}
Multiplication verifies the stated factorizations. The discriminant of
$G_c$ is $\delta(c)/36$, so when $\delta(c)\ne0$ the two critical
points $u,v$ are distinct, and
$h_c(u)=h_c(v)$ is equivalent to $P(c)(u-v)=0$, hence to $p(c)=0$.
For a completely explicit check of the double-factor identity, polynomial
division gives
\begin{multline*}
 h_c-Q(c)-G_c^2H_c\\
 =\frac{p(c)}{2592}
 \bigl(-5c^2-4c-30cz^2+36cz+24z^3-24z^2\bigr).
\end{multline*}
Finally division of $Q(c)-(32-52c-c^2)/1500$ by $p(c)$ has zero
remainder. These are polynomial identities over $\mathbb Q$, also
verified coefficient by coefficient in the exact certificate.
\end{proof}

\paragraph{B2.2 (Proved: all the Belyi functions, with no missing solutions).}
Let $c$ range over the three roots of $p$. Then
\begin{equation}\label{eq:B2belyi}
 \beta_c:\mathbb P^1_{\mathbb C}\longrightarrow\mathbb P^1_{\mathbb C},
 \qquad
 \beta_c(z)=\frac{1500z^3(z-1)^2(z-c)}{32-52c-c^2}
\end{equation}
gives exactly the three isomorphism classes of the required passport.
The roots of $G_c$ are its two double white vertices, and the roots of
$H_c$ are its two simple white vertices. No source automorphism of any
of these Belyi pairs is nontrivial.

\begin{proof}
The values of $p$ on $0,1,\ldots,6$ modulo $7$ are
$5,1,4,3,1,1,6$, respectively. A cubic over a field is reducible
if and only if it has a root, so $p$ is irreducible over $\mathbb Q$.
Its discriminant is $-2^8 3^9\ne0$, so it has three distinct roots.
Neither $0$ nor $1$ is a root, and irreducibility implies that none
is a zero of the nonzero quadratics $\delta$ or $32-52c-c^2$.
For any such $c$, the black divisor is therefore precisely
$3[0]+2[1]+[c]$ and the pole divisor is $6[\infty]$.
Equation (\ref{eq:B2factor}) gives
$\beta_c-1=G_c^2H_c/Q(c)$.
The residual quadratic $G_c$ has two simple roots. Its values at
$0,1,c$ are $c/2,(1-c)/3,c(c-1)/6$, so these roots are not black
vertices. Further,
\[
 \operatorname{disc}_z H_c=\frac{16+4c-11c^2}{9},\qquad
 \operatorname{Res}_z(G_c,H_c)
 =\frac{(c-1)(25c^3+18c^2-12c-4)}{36}.
\]
Modulo $p$, the latter equals
$(c-1)(5c^2+2c+2)/6$. Irreducibility again shows that both displayed
quantities are nonzero. Thus the white divisor has exactly the desired
multiplicities. The derivative has no further critical points; infinity
has ramification six. Hence (\ref{eq:B2belyi}) is a Belyi function with
the asserted passport.

Conversely, any pair of that passport has genus zero by
$3+4+1=6+2$, and the sole pole has order six. There is a unique source
coordinate sending its pole, its triple black vertex and its double
black vertex to $\infty,0,1$ respectively. In that coordinate its
function is $\lambda h_c$ with $\lambda\ne0$ and $c\ne0,1$.
Its two other critical points must be distinct and have common
value one. The derivative and remainder calculation force
$\delta(c)\ne0$, $p(c)=0$, and $\lambda Q(c)=1$.
Thus every pair occurs in our list. An isomorphism between two
normalized pairs must fix the uniquely characterized points
$\infty,0,1$, so it is the identity, and it then forces $c=c'$.
The same argument proves the automorphism assertion.
\end{proof}

\paragraph{B2.3 (Proved: the three combinatorial maps).}
With $\rho=(123)(45)(6)$, representatives are
\[
 \begin{array}{c|c|c}
 \text{map}&\alpha&\phi=\rho^{-1}\alpha\\\hline
 T_1&(14)(36)&(154362)\\
 T_0&(14)(56)&(156432)\\
 T_2&(14)(26)&(154326)
 \end{array}
\]
Unwritten points in $\alpha$ are fixed. The manuscript's tables I and II
on scan 106 are $T_1$ and $T_0$, respectively. The missing map is
the mirror $T_2$ of $T_1$. The map $T_0$ is its own mirror.
The function for the real root of $p$ is $T_0$, and the two nonreal
roots give $T_1,T_2$ in some order. We identify each nonreal map
unambiguously by the sign of the imaginary part of its normalized
parameter $c$; identifying it further with a chosen dart labelling is
unnecessary and is not claimed here.

\begin{proof}
Each table entry is checked by composition, and each face permutation
is a six-cycle, hence generates a transitive subgroup together with the
other permutations. The cycle counts give genus zero. To see
completeness without assuming the algebraic calculation, fix $\rho$.
Its centralizer has order $3\cdot2=6$. There are 45 involutions with
two two-cycles and two fixed points; the certificate lists all of them,
retains exactly 18 having $\rho^{-1}\alpha$ a six-cycle, and computes
their three disjoint centralizer orbits of size six. Equivalently, among
all $120\cdot45=5400$ ordered pairs of the required first two types,
there are 2160 admissible pairs and three simultaneous-conjugacy
orbits of size 720. The displayed representatives belong to different
orbits. This finite enumeration is exhaustive by its explicit loops,
and B2.2 supplies an independent completeness proof.
Orientation reversal replaces $\rho$ by $\rho^{-1}$ and leaves
$\alpha$ unchanged. Conjugation by $(23)$ restores our fixed $\rho$,
exchanges the first and third displayed involutions, and fixes the
second. The real root $c_0$ lies in $(1,3/2)$, since $p(1)<0<p(3/2)$
and $p$ has only one real root. The signs of $6G_c$ at $0,1,c$
place one critical point between $0$ and $1$, and the other between
$1$ and $c$. The real dessin therefore contains the path
$0-u-1-v-c$, with its other two white leaves attached to $0$;
this is $T_0$.
\end{proof}

\paragraph{B2.4 (Proved: fields of moduli, fields of definition and L5).}
Put $r_j=2^{1/3}\zeta_3^j$ for $j=0,1,2$, where $2^{1/3}$ is positive
real, and set
\[
 c_j=\frac{4+18r_j+6r_j^2}{25}.
\]
The field of moduli of $\beta_{c_j}$ relative to $\mathbb Q$ is
$\mathbb Q(c_j)=\mathbb Q(r_j)$, and (\ref{eq:B2belyi}) is a model
over that field. The three maps form one Galois orbit. Its permutation
group is $S_3$, with splitting field
$\mathbb Q(2^{1/3},\zeta_3)$ and quadratic subfield
$\mathbb Q(\sqrt{-3})$. Complex conjugation fixes $T_0$ and interchanges
the other two maps. Thus L5 is true when its ``cubic $P$'' means the
effective cubic $p$; it is false if it literally identifies the raw
remainder coefficient $P$ with a cubic.

\begin{proof}
Substitution and reduction by $r^3-2$ give $p((4+18r+6r^2)/25)=0$.
Irreducibility of $p$ proves that these are its three distinct roots.
Since $r_j^3-2$ is irreducible by Eisenstein at 2, both
$\mathbb Q(c_j)\subseteq\mathbb Q(r_j)$ have degree three and are
equal. For any $\sigma\in\operatorname{Gal}(\overline{\mathbb Q}/\mathbb Q)$,
coefficient conjugation sends $\beta_c$ to $\beta_{\sigma(c)}$.
By the normalization argument in B2.2, the two pairs are isomorphic
if and only if $\sigma(c)=c$. Their stabilizer is therefore exactly
$\operatorname{Gal}(\overline{\mathbb Q}/\mathbb Q(c))$, whose fixed
field is $\mathbb Q(c)$. This proves the moduli-field assertion directly,
without assuming that a field of moduli is always a field of definition.
The displayed coefficients prove descent to that field in this case.
Irreducibility gives transitivity. The nonsquare discriminant
$-2^8 3^9$ makes the transitive subgroup of $S_3$ the full $S_3$.
The description by cube roots identifies its normal closure and
quadratic subfield; alternatively the square root of the discriminant
gives the same quadratic field. The asserted complex-conjugation
action follows from the one real root and B2.3.
\end{proof}

We found this cubic already in Schneps, IV, Example 3, author-PDF
pages 27--28: writing her variable as $t$, her equation is
$25t^3-6t^2-6t-2=0$, which becomes our $p(c)=0$ upon $c=2t$.
We therefore make no priority claim for
the cubic or the orbit. Our proof above verifies the form actually
needed here instead of relying on a change of normalization left implicit.

\paragraph{B2.5 (Refuted and corrected: the collision and sign questions).}
The two roots of $\delta$ are
$16/25\pm(12/25)i$, not $16/25\pm(24/25)i$.
Both are zeros of the raw $P$ on scan 109, answering the question
on scan 112 affirmatively. They give the different passport
$[3,2,1;3,1,1,1;6]$ after normalization, so they must be excluded
from B2.2. The undetermined sign in
$(\alpha-2)(\alpha-c)+\beta+1$ on scan 108 is minus in
$-(5c^2-10c-4)/36$, and the numerator on scan 109 is $f-1$.
The divisor on scan 111 is $(z^2-\alpha z-\beta)^2$.
The other doubtful coefficient on scan 108 satisfies
$\alpha(\alpha^2+2\beta)=(5c+4)(25c^2+4c+16)/216$.
For the final rational-function calculation on scan 112, if
$f=\lambda\varphi/\psi$ with $\lambda\ne0$ and coprime monic
$\varphi=\prod_i(z-a_i)^{p_i}$ and
$\psi=\prod_j(z-b_j)^{q_j}$, the identity in $\mathbb C(z)$ is
\[
 \frac{f'}f=\sum_i\frac{p_i}{z-a_i}-\sum_j\frac{q_j}{z-b_j}.
\]

\begin{proof}
The quadratic formula gives
$\operatorname{disc}\delta=-576$ and the stated roots. The
factorization in B2.1 shows $\delta\mid P$.
For a root $c$ of $\delta$, the repeated residual critical point is
$u=(5c+4)/12$. The evaluations of $G_c$ at $0,1,c$ used in B2.2
show that $u$ is none of these three points. Also,
\[
 Q(c)\equiv-\frac{2(139c-82)}{9375}\pmod{\delta(c)}
\]
is nonzero at either nonreal root. Thus $h_c/Q(c)$ is defined,
and its derivative has order exactly two at $u$, giving ramification
index three over one. There are no other residual critical points,
so the other three points over one are simple. This proves the altered
passport and explains why the raw quintic cannot be used without
the discriminant exclusion. Direct substitution of
$\alpha=(5c+4)/6$ and $\beta=-c/2$ gives the asserted sign on scan
108. The equation $f(u)=f(v)=1$ forces the numerator $f-1$;
the monic polynomial of $u,v$ is $z^2-\alpha z-\beta$.
The remaining coefficient follows by substituting $\alpha,\beta$;
the logarithmic derivative identity follows from the product and
quotient rules and fixes the last minus sign on scan 112.
All alternatives marked by doubt in these mathematical expressions
are thereby resolved by polynomial identities; no appeal to uncertain
handwriting is needed for the resulting theorems.
\end{proof}

\section{B.3. Holomorphic automorphisms preserving every cell}

Here a clean holomorphic map means the data on 138:129:
a nonconstant clean Belyi map
$\beta:X\longrightarrow\mathbb P^1_{\mathbb C}$ on a connected compact
Riemann surface, with graph $K=\beta^{-1}([0,1])$ and vertex set
$S=\beta^{-1}(0)$. White points are interior edge marks, not vertices
of the ordinary map. Let $G_{\mathrm{cell}}(X,K,S)$ denote the group of
biholomorphisms of $X$ fixing every vertex, carrying each edge to
itself, and carrying each face to itself. We do not require them to
preserve $\beta$. Let $H$ be the subgroup inducing the identity on
every dart. A loop has two darts at its vertex, so fixing it setwise
need not fix both darts.

\paragraph{B3.1 (Proved: the complete compact exception list).}
Up to biholomorphic equivalence of the embedded map, the nontrivial
groups $G_{\mathrm{cell}}$ are exactly the following:
\begin{enumerate}
\item The sphere with one vertex and one loop (case I): in a coordinate
with $K=\mathbb R\cup\{\infty\}$ and $S=\{\infty\}$,
$G_{\mathrm{cell}}=\{z\mapsto az+b:a>0,\ b\in\mathbb R\}$.
\item The sphere with two vertices and $n\geq1$ edges joining them
(the orange of order $n$): in a coordinate with vertices $0,\infty$
and $n$ equally spaced rays as edges,
$G_{\mathrm{cell}}=\{z\mapsto az:a>0\}$.
Case II is $n=2$, and case III is $n=1$.
\item The genus-$g$ one-vertex one-face map of B1.3, for $g\geq1$:
$G_{\mathrm{cell}}\simeq C_2$, and its nonidentity element reverses
the two darts of every edge.
\end{enumerate}
Every other connected finite clean holomorphic map has trivial
$G_{\mathrm{cell}}$. The subgroup $H$ equals the whole displayed
group in the two sphere families and is trivial in positive genus.
For a disconnected finite map the answer is the product of these
groups over its connected components; cellwise preservation permits
no permutation of components.

\begin{proof}
Every holomorphic cellwise automorphism preserves orientation and
therefore induces an element of $K_C$ from B1.3. Every element of
$K_C$ is realized by a unique Belyi automorphism. To spell out this
last assertion, restrict $\beta$ to its unbranched covering over
$U=\mathbb P^1\setminus\{0,1,\infty\}$. A permutation of one fibre
commuting with monodromy defines a covering transformation by
path lifting: the commuting condition makes the result independent
of the chosen path. On each evenly covered open disk it is the
composition of inverse branches, hence holomorphic. At a point over
$0,1$ or $\infty$, local coordinates identify the covering with
$z\mapsto z^e$, and its lift extends across the centre. The extension
is unique by continuity. This proves the assertion with source and
target $\operatorname{Aut}(C)\longrightarrow\operatorname{Aut}(X,\beta)$.
Consequently
$G_{\mathrm{cell}}\longrightarrow K_C$ is a split surjection with
kernel $H$.

For $h\in H$, choose a vertex $s$. The graph is locally a finite
union of analytic rays with distinct tangent directions: take a
local coordinate in which $\beta=w^d$. Since $h$ fixes each outgoing
dart, its complex derivative at $s$ is a positive real number.
If $X$ has positive genus, its group of holomorphic automorphisms
fixing $s$ is finite. We use Wendl, Corollary 4.2.7, applied to the
stable pointed compact surface $(X,\{s\})$; the hypotheses hold for
every genus at least one \cite{Wendl2015B}. Thus $h$ has finite order,
and positivity forces its derivative to equal one. A finite-order
holomorphic germ with derivative one is the identity: in a local
coordinate $w$ and for order $m$, the function
$m^{-1}\sum_{j=0}^{m-1}w\circ h^j$ is a coordinate with nonzero
derivative and is fixed by $h$. The identity theorem then gives
$h=1$ on $X$. Hence $H=1$ in positive genus, and B1.3 gives exactly
the third family and its stated group.

Suppose now that $X=\mathbb P^1$. By B1.3, $K_C=1$ in genus zero,
so every cellwise automorphism belongs to $H$. If there are at least
three vertices, a M\"obius transformation fixing all of them is the
identity. If there are exactly two, place them at $0,\infty$.
Every automorphism fixing them is $h(z)=az$, and fixing an outgoing
dart makes $a>0$. Assume $a\ne1$, replacing $h$ by $h^{-1}$ if
necessary so $a>1$. The closure of every edge is invariant and
compact. For any interior point $z$ of the edge, the sequences
$a^kz$ tend to $0$ and $\infty$ as $k$ tends to opposite infinities.
Thus both vertices lie on that edge, so every edge joins them.
Its tangent direction at $0$ is the limit of the arguments of
$a^{-k}z$, which is just $\arg z$. Hence all its points have that
same argument and the edge is a ray. The local form $\beta=w^n$
forces these $n$ tangent rays to be equally spaced. Conversely all
positive dilations preserve each ray and each intervening sector.
This proves the second family and that its full group is precisely
$\mathbb R_{>0}$.

It remains to consider a single vertex, which we place at $\infty$.
Every automorphism fixing it has the form $h(z)=az+b$.
The derivative condition at infinity gives $a>0$. If $a\ne1$, a
translation puts the other fixed point at zero. For any nonfixed
point in an edge, its two orbit limits are $0$ and $\infty$.
The point $0$ must therefore belong to the closure of every edge.
It is not a vertex, so it belongs to an edge interior; distinct
edge interiors are disjoint. There is consequently only one edge.
If $a=1$ and $b\ne0$, the sequences $z+kb$ tend to $\infty$ with
the two tangent directions $\pm1/b$ in the coordinate $w=1/z$.
Every invariant loop must have exactly these two tangent directions
there. Distinct darts at a clean Belyi vertex have distinct tangent
directions, so again there can be only one edge. Thus any nonidentity
element forces the one-loop map. A clean Belyi realization of this
map is $1/(1+z^2)$, with graph $\mathbb R\cup\{\infty\}$ and sole
vertex $\infty$; uniqueness of its abstract Belyi pair, or direct
normalization of a degree-two rational function, identifies every
realization with this one. Its M\"obius automorphisms fixing infinity
are affine; preserving the real circle forces real coefficients,
and preserving the two faces separately forces $a>0$. This is
exactly the first group in the list.
\end{proof}

\paragraph{B3.2 (Proved: explicit models and the manuscript corrections).}
The orange of order $n$ has clean Belyi function
\[
 b_n(z)=-\frac{4z^n}{(z^n-1)^2}.
\]
Its graph consists of the $n$ rays $z^n\in\mathbb R_{\leq0}$,
with vertices $0,\infty$. In particular every positive dilation is
a nonidentity cellwise holomorphic automorphism when its scale is
not one, for every $n\geq1$. The assertion on 138:131 that there
are none for $n\geq2$ is false. The full groups in I, II and III
are the real groups stated in B3.1, resolving either real/complex
reading of the doubtful symbols on scan 130.
The third family has explicit model
\[
 X_g:\ y^2=x^{2g+1}-x,\qquad \beta_g(x,y)=x^{2g},
 \qquad (x,y)\longmapsto(x,-y)
\]
on the smooth projective completion of the indicated affine curve.
The displayed involution generates its cellwise group.

\begin{proof}
For $t=z^n$ and $0<s\leq1$, the equation
$-4t/(t-1)^2=s$ has two real nonpositive roots (a double root
$t=-1$ at $s=1$), since its discriminant is $16(1-s)$.
At $s=0$ its two projective roots are $0,\infty$.
This proves the graph description. The zeros at $0,\infty$
have order $n$, all roots of $z^n=1$ are double poles, and
$b_n-1=-(z^n+1)^2/(z^n-1)^2$, showing cleanliness and all
branch values. Positive scaling preserves each ray and each sector;
for scale different from one it is not the identity.

The polynomial $x^{2g+1}-x$ has distinct roots in characteristic
zero, and the double cover $X_g\longrightarrow\mathbb P^1_x$ is
branched at those $2g+1$ roots and at infinity. Riemann--Hurwitz
gives genus $g$. The function $x^{2g}$ on $X_g$ has degree $4g$,
a unique zero of order $4g$ at $x=0$, a unique pole of order $4g$,
and $2g$ points over one, each of index two. There are no other
critical values. For $\zeta=\zeta_{4g}$, the transformation
$(x,y)\mapsto(\zeta^2x,\zeta y)$ is a deck transformation of order
$4g$, so the monodromy is regular cyclic. A generator around the
unique black vertex is a full $4g$-cycle, and the order-two
involution over the white vertices is its $2g$-th power. Thus the
dart datum is precisely B1.3, and that power is $(x,y)\mapsto(x,-y)$.
B3.1 proves that the full cellwise group is the stated $C_2$.
The finite genus-one counterexample has been exhaustively verified
in the certificate, and the displayed permutation models have also
been checked for genera 1 through 8; the proof covers every $g$.
\end{proof}

\paragraph{B3.3 (Proved: uniqueness and the distinction from isotopy).}
On a fixed compact $(X,K,S)$, the clean Belyi functions with these
same embedded cells are parametrized freely and transitively by
$H$ through precomposition. Nonuniqueness occurs exactly for case I
and for the oranges of every order. The finite $C_2$ in positive
genus preserves the Belyi function itself and creates no additional
function. Cellwise preservation is therefore not equivalent to
isotopy to the identity, contrary to the identification on 138:129.
Specifying every face centre kills the continuous groups in I and
all oranges; specifying every white point kills them in the oranges
but leaves the positive dilation subgroup in I.

\begin{proof}
Given two such Belyi functions, the identity of their abstract
oriented maps gives the isomorphism of their unbranched coverings
with the prescribed identification of darts. By the continuation
argument in B3.1 it extends to a biholomorphism of $X$, lying in
$H$, and intertwines the functions. Conversely any $h\in H$
preserves exactly the required embedded cells. Its precomposition
acts trivially on a function only if it is a Belyi automorphism;
the latter acts freely on darts, so $h=1$. This proves the
parametrization and the uniqueness criterion.
For a concrete failure of the isotopy identification, take $g=1$
in B3.2. The involution reverses both loops in the one-vertex
one-face cell structure, hence induces $-1$ on its two-dimensional
first integral homology. The cellular boundary of its sole two-cell
is zero, so these two oriented loop classes form a basis. An
automorphism isotopic to the identity induces the identity on
homology, and this involution does not.
In case I the face centres are $i,-i$ and the white point is $0$:
$z\mapsto az+b$ with $a>0,b\in\mathbb R$ fixes both centres
only for $a=1,b=0$, whereas fixing the white point only requires
$b=0$. For an orange, every centre and every white point is a
nonzero finite number, so a positive dilation fixing any one of
them has scale one. These direct calculations prove the final
assertions.
\end{proof}

\paragraph{B3.4 (Proved under stated hypotheses: locally finite maps).}
The same conclusion extends to connected locally finite clean
holomorphic cellular maps with compact edge closures and ordinary
disk faces attached along their edge cycles. On a noncompact
Riemann surface every cellwise holomorphic automorphism is the
identity. These are the additional hypotheses of this infinite-map
extension; arbitrary holomorphic maps to the sphere with asymptotic
tracts do not automatically define such a cellular map.

\begin{proof}
The dart-kernel argument of B1.3 works without finiteness until it
shows that a nontrivial kernel forces one $\rho$-cycle. Local
finiteness makes the valence at that sole vertex finite, so all
darts are then finite in number. There are finitely many edge and
face cycles, and their ordinary compact cell attachments make
the resulting surface compact. Thus on the noncompact surface
the combinatorial cellwise kernel is trivial, and a cellwise
automorphism $h$ fixes every dart.
By uniformization (Wendl, Theorem 4.2.1 and the immediately following
quotient classification, opened), the universal cover is the disk
or the plane. In the disk case, lift $h$ so that it fixes a chosen
lift of a vertex. The lift is a disk automorphism fixing that point;
after putting that point at zero it is a rotation by the Schwarz
lemma. Its derivative is positive real because it fixes a dart,
so the rotation is the identity, and $h=1$.
In the plane case, the noncompact surface is $\mathbb C$ or
$\mathbb C^*$. A holomorphic automorphism of $\mathbb C^*$ fixing
a point is either the identity or a conjugate of inversion, whose
derivative at its fixed point is $-1$; the latter cannot fix a dart.
On $\mathbb C$, an automorphism fixing a vertex is an affine
dilation about it, with positive real multiplier. If the multiplier
differs from one, the two-sided orbit of any other point is
unbounded in one direction. Such an orbit cannot be contained in
an invariant compact edge closure. Hence the multiplier is one
and the automorphism is the identity. This exhausts all cases.
\end{proof}

\paragraph{Verification and provenance for B.1--B.3.}
The exact program \path{code/B1_B2_B3_certificate.py} and its
saved output verify every finite counterexample used above, the
complete $S_6$ enumeration, polynomial identities, discriminants,
excluded degeneracies and cube-root parametrization. Its scope
is explicitly finite; general statements use the proofs above.
Literature searches included the passport, the sextic family,
cell-preserving holomorphic map automorphisms, and the centralizer
description. The sextic cubic is already in the opened Schneps
source. We found no source for the precise combined cellwise list
B3.1 during these searches and assert no priority for it.
